% Separate actual local calculation, after S16.
% GPT-6 Astra Pro, /root/fixed_surface, 2026-11-22 UTC; effort not exposed.

\section{HD-F1-S17: normality and the actual tangent cones of the norm image}
\label{sec:f1-norm-fourfold-normality}

For the actual covers and objects of HD-F1-S16, the Abel image $Y$ is
normal with rational singularities. Thus the actual candidate is
\[
                        Ra_*\mathcal O_T=\mathcal O_Y.
 \tag{F1S17.1}
\]
At every image $y$ of a complete-pencil fibre $E=\mathbb P^1$, write
\[
                  N^*_{E/T}=\bigoplus_{i=1}^3\mathcal O_E(b_i).
\]
Then $b_i\geq0$, $\sum b_i=3$, and its actual Zariski tangent cone is
\[
 \operatorname{Spec}\left(
       \bigoplus_{m\geq0}H^0\left(\mathbb P^1,
                \operatorname{Sym}^m\!\left(\bigoplus_i\mathcal O(b_i)\right)
                                     \right)\right).
 \tag{F1S17.2}
\]
In particular $\dim T_yY=6$, the Hilbert--Samuel multiplicity is three,
and the tangent-cone Hilbert series is $(1+2t)/(1-t)^4$. Up to order,
the three possible splitting types are $(1,1,1),(0,1,2),(0,0,3)$.
No assertion choosing among them, counting the pencils, or identifying
the completed local ring with its tangent cone is made.

\paragraph{1. The exact projected Petri monad along a pencil.}
Fix the actual line bundle $L$ defining $E$, put
$U=H^0(Z,L)$ (dimension two), $A_L=H^0(Z,K_ZL^{-1})$ (dimension six),
and put
\[
             V=H^0(Z,K_Z)/\pi^*H^0(C,K_C),\qquad\dim V=8.
\]
The norm identity gives the actual isomorphism
$K_Z\simeq L\otimes\sigma^*L\otimes\sigma^{2*}L$.
For $[u]\in\mathbb P U$, define the nonzero quadratic section
$\alpha(u)=\sigma^*u\,\sigma^{2*}u\in A_L$.
Multiplication by $u$, followed by projection to $V$, defines a bundle
map $A_L\otimes\mathcal O_E(-1)\to V\otimes\mathcal O_E$.
Its kernel at $[u]$ consists precisely of the forms whose product with
$u$ descends to $C$ and vanishes along $D=(u)$. The exact maximum-
multiplicity evaluation of S13 makes this intersection one dimensional.
It is the line spanned by $\alpha(u)$, since
$u\alpha(u)=\pi^*\operatorname{Nm}(u)$ is the canonical norm section.
The quadratic map $\alpha$ is nowhere zero, since a product of two
nonzero sections of line bundles on the integral curve is nonzero.
Thus the actual kernel line bundle is $\mathcal O_E(-3)$.

The dual of the normal map $N_{E/T}\to\mathcal O_E^{\oplus8}$ in S16
has precisely this image as kernel, fibre by fibre, and all ranks are
constant. We have proved the exact sequence
\[
 0\longrightarrow\mathcal O_E(-3)
 \xrightarrow{\alpha} A_L\otimes\mathcal O_E(-1)
 \longrightarrow V\otimes\mathcal O_E
 \longrightarrow N^*_{E/T}\longrightarrow0.
 \tag{F1S17.3}
\]
Let $\mathcal K$ denote the image of its second arrow. The first short
exact sequence gives
$H^0(\mathcal K)=H^1(\mathcal O_E(-3))\simeq\mathbb C^2$ and
$H^1(\mathcal K)=0$, because all cohomology of $\mathcal O_E(-1)$
vanishes. The second short exact sequence therefore gives the actual
surjection
\[
                    V\twoheadrightarrow H^0(E,N^*_{E/T}),
              \qquad \dim H^0(E,N^*_{E/T})=6.
 \tag{F1S17.4}
\]
This is the map induced by the ambient cotangent coordinates at $y$;
it is not an abstract isomorphism of two spaces of the same dimension.

For clarity, dualizing (F1S17.3) gives
$0\to N_{E/T}\to\mathcal O_E^8\to\mathcal O_E(1)^6\to
\mathcal O_E(3)\to0$. Even though the last map is surjective on global
sections, this does not make the intervening map between two
8-dimensional global-section spaces an isomorphism. Thus (F1S17.3)
does not prove $H^0(N_{E/T})=0$ or isolate the pencil. Those conclusions
are not used here.

\paragraph{2. Exact formal-functions filtration of the normalization ring.}
Let $R=\mathcal O_{Y,y}$, and let $B=(a_*\mathcal O_T)_y$, a finite
normalization algebra over $R$ by S16. Since the scheme fibre of $a$
is connected, the finite normalization fibre over $y$ has one point;
otherwise its disjoint inverse images would disconnect the fibre of
$a$. Let $\widehat B$ be its completion over $\mathfrak m_R$.
By formal functions in degree zero,
\[
       \widehat B=\varprojlim_m H^0(E_m,\mathcal O_{E_m}),
 \quad E_m\text{ defined by }\mathcal I^{m+1},
 \quad\mathcal I=\mathfrak m_R\mathcal O_T.
 \tag{F1S17.5}
\]
Define $F^m\widehat B$ as the kernel of restriction to $E_{m-1}$,
with $F^0=\widehat B$. The symmetric-conormal higher-cohomology
vanishing in S16 makes every restriction on these thickenings
surjective and gives multiplicative identifications
\[
 \operatorname{gr}_F^m\widehat B
                    =H^0(E,\operatorname{Sym}^mN^*_{E/T}).
 \tag{F1S17.6}
\]
The filtration is complete and separated; its topology is the
$\mathfrak m_R$-adic topology by the topological statement in formal
functions. Each right side is computed using the actual regular-
embedding isomorphism $\mathcal I^m/\mathcal I^{m+1}=
\operatorname{Sym}^mN^*_{E/T}$, so (F1S17.6) respects multiplication.

The graded algebra on the right is generated by degree one. Indeed,
expand the symmetric power using the splitting with $b_i\geq0$.
Every summand is $H^0(\mathcal O(\sum b_i k_i))$ with
$\sum k_i=m$. Multiplication of binary homogeneous polynomials gives
a surjection from the corresponding products of
$H^0(\mathcal O(b_i))$. Explicitly every degree-$\sum b_i k_i$
monomial in two variables can be distributed among factors of the
prescribed nonnegative degrees. This proves the required surjectivity
without imposing any of the three possible splitting types.

\paragraph{3. Ambient coordinates prove normality, then identify the tangent cone.}
Choose local analytic coordinates $x_1,\ldots,x_8$ on the smooth
ambient $P$ at $y$. Their pullbacks belong to $F^1\widehat B$, and
their classes in $F^1/F^2$ give exactly (F1S17.4). Consequently these
eight classes generate the associated graded algebra. Successive
lifting in each filtration degree expresses every element of
$\widehat B$ as a convergent formal power series in the pullbacks of
the $x_i$: at each step subtract a homogeneous polynomial with the
required associated-graded class; completeness gives the limiting
series. Hence the actual map
\[
                  \mathbb C[[x_1,\ldots,x_8]]
                               \longrightarrow\widehat B
\]
is surjective. It factors through $\widehat R$, because it is the
completion of the given morphism to the actual reduced image $Y$.
The injective finite map $R\to B$ remains injective under completion;
its completed map is now also surjective. Faithfulness of completion
on finite modules gives $R=B$. Thus the normalization is an isomorphism
at $y$. Away from complete pencils the Abel map has one-point fibres
and injective differential, so it is locally an embedding and the
image is smooth there. This proves the global normality assertion.
Its rational singularities and (F1S17.1) follow from S16.

The maximal ideal of $\widehat B$ is $F^1$, since its residue ring is
$H^0(E,\mathcal O_E)=\mathbb C$. Generation of (F1S17.6) in degree one
shows $(F^1)^m+F^{m+1}=F^m$. Iterating and using the equality of the
complete noetherian topologies gives $(F^1)^m=F^m$: the ideal
$(F^1)^m$ is closed, so the remainders in $F^N$ tend to zero inside it.
Thus (F1S17.6) is the actual maximal-ideal associated graded ring, which
proves (F1S17.2).

Finally, for $m\geq0$, the number of triples $k_i\geq0$ with sum $m$
is $\binom{m+2}{2}$. Symmetry gives
$\sum_{k_1+k_2+k_3=m}\sum_i b_ik_i
  =m\binom{m+2}{2}$, because $\sum b_i=3$. Therefore
\[
 \dim\operatorname{gr}_{\mathfrak m}^m\mathcal O_{Y,y}
          =(m+1)\binom{m+2}{2},\qquad
 \sum_{m\geq0}(m+1)\binom{m+2}{2}t^m
                         =\frac{1+2t}{(1-t)^4}.
 \tag{F1S17.7}
\]
This gives embedding dimension six and multiplicity three. For
$(b_i)=(1,1,1)$ the ring in (F1S17.2) is explicitly the rank-one
$2\times3$ determinantal ring. For $(0,1,2)$ it has one independent
polynomial variable and the three minors of
$\left(\begin{smallmatrix}x_0&y_0&y_1\\x_1&y_1&y_2\end{smallmatrix}\right)$;
for $(0,0,3)$ it has two independent polynomial variables and the three
minors of
$\left(\begin{smallmatrix}y_0&y_1&y_2\\y_1&y_2&y_3\end{smallmatrix}\right)$.
These presentations follow from the displayed multiplication of binary
forms. They describe the tangent cone, while the full local analytic
relations and the Hochschild trace are still separate actual
calculations.
